% Diagrama Mermaid en hoja carta apaisada.
% El trazo viene de niveles-concrecion-mermaid.mmd, renderizado a PNG a 322 dpi:
%   npx @mermaid-js/mermaid-cli -i niveles-concrecion-mermaid.mmd \
%       -o niveles-concrecion-mermaid.png -b white -s 4
% Compilar:  tectonic docs/actividad/niveles-concrecion-diagrama.tex

\documentclass[11pt]{article}

\usepackage[T1]{fontenc}
\usepackage[utf8]{inputenc}
\usepackage[spanish,es-noindentfirst]{babel}
\usepackage{lmodern}
\usepackage[letterpaper,landscape,top=0.9cm,bottom=0.8cm,left=1cm,right=1cm]{geometry}
\usepackage{graphicx}
\usepackage{xcolor}

\pagestyle{empty}
\setlength{\parindent}{0pt}

\definecolor{tinta}{HTML}{1B2130}
\definecolor{gris}{HTML}{4A5265}

\newcommand{\sep}{\;\textbullet\;}

\begin{document}

\begin{center}
  {\Large\bfseries\color{tinta} Niveles de concreción de la gestión educativa}\\[3pt]
  {\small\color{gris} Diagrama de flujo \sep Modelo de Gestión Educativa Estratégica (SEP)
   \sep Temas selectos de Pedagogía}
\end{center}

\vfill
\begin{center}
  \includegraphics[width=\textwidth,height=0.80\textheight,keepaspectratio]%
    {niveles-concrecion-mermaid.png}
\end{center}
\vfill

{\scriptsize\color{gris}\textbf{Referencia:}\ \ Vázquez Herrera, E. (Coord.). (2010).
\textit{Modelo de gestión educativa estratégica} (2.\textsuperscript{a}~ed.). Secretaría de
Educación Pública. \sep Elaboración propia a partir de las pp.~55--64
(numeración del libro impreso).}

\end{document}
